%% -*- coding: utf-8 -*-
\documentclass[12pt,pagesize,paper=192mm:108mm,landscape]{scrbook} 
%1920x1080 1280x720
\areaset[current]{192mm}{108mm}
\usepackage{calc}
\usepackage[T2A]{fontenc}
\usepackage[utf8]{inputenc}
\usepackage[english,russian]{babel}
\usepackage{microtype}
\usepackage{misccorr}
\usepackage{cmap}
%\usepackage[unicode=true]{hyperref}
\usepackage{graphicx}
\usepackage{amssymb}
\usepackage{amsmath}
%\usepackage{srcltx}
\usepackage{textcomp}
\usepackage{xspace}
%научные символы и смайлики \smiley \frownie
\usepackage{wasysym}
\usepackage{ccicons}
\begin{document}
\begin{titlepage}
  \vspace*{-0.5em}
  \begin{center}    
    \hspace*{3em}
    \begin{minipage}[t]{1.5em}
      \includegraphics[width=\textwidth]{../BINP-logo}
    \end{minipage}\hfill
    \begin{minipage}{0.115\linewidth}
    \includegraphics[width=\textwidth]{../NSU-logo}
    \end{minipage}
    \hfill
    \hspace*{4.5em}

    Кафедра теоретической физики физического факультета НГУ

    \Large
    Профессор Грабовский А.\,В.

    \huge
    \textbf{Общая теория относительности}

    \Large
    Семинар № 10
    \vfill

    \normalsize
    \begin{minipage}{0.83\linewidth}
      Оценка необходимости применения ОТО для описания
      вселенной. Оценка расстояния до ближайшей звезды при условии,
      что поток световой энергии от нее, который мы видим, в 10--11 раз
      меньше потока энергии от Солнца. Оценка гравитационного красного
      смещения света, вылетающего из типичной галактики, сравнение
      этого красного смещения с хаббловским. Оценка относительной
      плотности энергии звездного света, считая, что звезды светили
      все время жизни вселенной с той же плотностью светимости, что и
      сейчас. Оценка изменения длины плеча ($\sim10$\,км) гравитационного
      интерферометра при прохождении гравитационной волны из кластера
      Девы, вызванной изменением гравитационного поля звезды солнечной
      массы.
    \end{minipage}
    \vfill

    \normalsize \ccbysa\hspace{0.5em}  Новосибирск 2022
  \end{center}
\end{titlepage}
\end{document}
